\documentclass[11pt,a4paper]{article}
\usepackage[margin=19mm,headheight=15pt]{geometry}
\usepackage{fontspec}
\setmainfont{Noto Serif}
\usepackage{amsmath,amssymb,mathtools}
\usepackage{xcolor}
\usepackage{fancyhdr}
\usepackage{enumitem}
\usepackage{microtype}
\setlength{\parindent}{0pt}
\setlength{\parskip}{6pt}
\definecolor{navy}{HTML}{183153}
\definecolor{lightgray}{HTML}{F3F5F7}
\pagestyle{fancy}\fancyhf{}
\lhead{\small ME 5309 -- Mid-Semester Examination}
\rhead{\small ME 5309}
\cfoot{\thepage}
\begin{document}
\begin{center}
{\Large\bfseries ME 5309 -- Mid-Semester Examination}\\[3pt]
{\large Answer Booklet}\\[8pt]
\end{center}
\textbf{Instructions:} Attempt any three questions. Each question carries 10 marks. Show working where appropriate.
\section*{Question 1 \hfill [10 marks]}
\subsection*{Q1(a) \hfill [4 marks]}

(i) f\_x = 16/5 = 3.2 pixels, using focal length 16 mm and pixel size 5 microm.
(ii) Sensor width is 6.4 mm, so field width = 400x6.4/16 = 160 mm and scale = 0.125 mm/pixel.

\subsection*{Q1(b) \hfill [3 marks]}

ADC step = 5/4096 = 1.2207 mV.
I divide by 0.1 mV/N, giving ΔF = 12.207 N.

\subsection*{Q1(c) \hfill [3 marks]}

\[\begin{aligned}F = 0.8187, G = 0.0906\end{aligned}\]

\section*{Question 2 \hfill [10 marks]}
\subsection*{Q2(a) \hfill [2 marks]}

The state model is ẋ₁ = x₂, ẋ₂ = (τ-0.05x₂-2.4525 sin x₁)/0.02.
Equilibrium requires τ₀ = 2.4525 sin θ\_e, so τ₀ must be between -2.4525 and +2.4525 N m.

\subsection*{Q2(b) \hfill [2 marks]}

\[\begin{aligned}A = [[0,1], [-61.3125,-2.5]], B = [[0],[50]]\end{aligned}\]

\subsection*{Q2(c) \hfill [4 marks]}

At steady state θ̇ = 0, and I assume θss = θd = 90° because the PD controller tracks the desired angle.
Thus the steady-state error is zero.

\subsection*{Q2(d) \hfill [2 marks]}

The gravity term changes the equilibrium, so the steady-state error is still nonzero.
The poles should change.

\section*{Question 3 \hfill [10 marks]}
\subsection*{Q3(a) \hfill [2 marks]}

Finite differencing amplifies encoder noise, making the velocity estimate noisy.

\subsection*{Q3(b) \hfill [2 marks]}

For H(s)=ωc/(s+ωc), ωc=2π(20)=125.664 rad/s and T=0.001 s.
Tustin gives H(z)=α(1+z⁻¹)/[(1+α)+(α-1)z⁻¹], α=ωcT/2=0.0628319.
Thus:\[\begin{aligned}y[k]=0.881765y[k-1]+0.0591175(x_k+x[k-1])\end{aligned}\]
DC gain = (2x0.0591175)/(1-0.881765) approximately 1.000.\par\smallskip $\omega_c=2\pi(20)=125.664$ rad/s. With $T=0.001$ s, $\alpha=\omega_cT/2=125.664(0.001)/2=0.062832$. The recurrence coefficients are approximately $0.881765$ on the previous output and $0.0591175$ on each input sample. Check: $2(0.0591175)/(1-0.881765)\approx1$, so the DC gain is close to one.

\subsection*{Q3(c) \hfill [4 marks]}

The characteristic polynomial is s²+(ℓ₁+2)s+(2ℓ₁+ℓ₂).
I match it to s²+10s+100 and get ℓ₁=8, ℓ₂=84.

\subsection*{Q3(d) \hfill [2 marks]}

An observer can estimate velocity from position measurements.
A filter can be safer because it smooths noise.

\end{document}